% !TEX root = ../main.tex
\section{Assumptions and Justifications}
\label{sec:assumptions}

\subsection{Model Basis}

埃博拉患者通常在出现症状后才具有传染性，病亡者的遗体仍可能造成感染，因此仅用易感、感染和康复三种状态（经典SIR模型）不足以区分疾病进程与救治情况\cite{whofactsheet}。Legrand 等人建立的 SEIHFR 模型将人群分为易感者（$S$）、潜伏者（$E$）、社区发病者（$I$）、住院患者（$H$）、死亡但尚未安葬者（$F$）以及已退出传播过程者（$R$），分别考虑社区、医院和葬礼中的传播\cite{legrand2007}。图~\ref{fig:seihfr-basis} 给出这一状态划分。

\begin{figure}[htbp]
\centering
\begin{tikzpicture}[
  state/.style={draw,rounded corners=2pt,minimum width=12mm,minimum height=9mm},
  flow/.style={-{Latex[length=1.8mm]},thick}
]
\node[state] (s) at (0,0) {$S$};
\node[state] (e) at (2,0) {$E$};
\node[state] (i) at (4,0) {$I$};
\node[state] (h) at (6,0) {$H$};
\node[state] (f) at (6,-1.6) {$F$};
\node[state] (r) at (8,-1.6) {$R$};
\draw[flow] (s) -- (e);
\draw[flow] (e) -- (i);
\draw[flow] (i) -- (h);
\draw[flow] (i) -- (f);
\draw[flow] (h) -- (f);
\draw[flow] (f) -- (r);
\draw[flow] (h) -- (r);
\draw[flow] (i.north) -- ++(0,0.8) -| (r.north);
\end{tikzpicture}
\caption{基础 SEIHFR 模型的状态转移示意，依据文献\cite{legrand2007}重绘。$R$ 包含康复者与已安葬的死者；箭头表示状态转移，不表示传播途径。}
\label{fig:seihfr-basis}
\end{figure}

该模型采用人群均匀混合的近似，没有逐一描述家庭成员及家属尸体、医护人员和其他接触者之间的关系。Legrand 等人也指出，均匀混合可能不足以反映家庭接触对传播的影响\cite{legrand2007}。这属于模型对接触结构的简化，不能等同于数值模拟中的随机采样错误。对于本题，还需要区分患者所在地区及地区之间的人口流动，单一人群的汇总模型无法直接给出这些信息。

本文保留六种基础状态，以 $j$ 表示所在省份、$k$ 表示家庭型，记相应人数为 $X_{jk}(t)$，其中 $X\in\{S,E,I,H,F,R\}$。省份分为八组，家庭型按城乡分为两组。这里的“家庭型”是人口类别，不是一个具体家庭，也不代表已经获得了逐户的亲属关系数据。第二个索引用于保留组间差异，组内仍作平均化处理。接种与治疗情景另设 $V$、$T$ 两种状态；它们不是所有基线计算都必须包含的舱室。各状态的完整定义与方程组在后文模型建立部分给出。

\subsection{Assumptions}

\begin{enumerate}[label=\textbf{A\arabic*.},leftmargin=2.6em,itemsep=0.6em]

\item \textbf{人口分类互不重叠。}
同一时刻，每个人只计入一个省份、一种家庭型和一种疾病状态。短期出行只改变所在省份，不改变家庭型。这样可以避免在省际汇总时重复计数；研究期内发生的长期迁居或家庭型变化暂不单独处理。

\item \textbf{不计自然出生和非疫病死亡。}
模型关注本轮疫情及随后一段时间的变化，暂不考虑与疫情无关的人口增减。省际流动只改变各省的人口分布，不凭空增加或减少模型中的人数。疫情死亡仍须保留在状态转移和人口核算中，不能因忽略自然死亡而一并省略。

\item \textbf{潜伏者不参与传播，潜伏期结束后进入发病状态。}
这一处理与 WHO 关于症状出现前不传播疾病的说明一致\cite{whofactsheet}，也沿用了基础 SEIHFR 的潜伏状态定义。模型不另设无症状感染舱，因此不用于估计无症状感染的比例。

\item \textbf{康复者与已安葬的死者退出本模型描述的传播过程。}
二者均计入 $R$，但并不具有相同的医学含义。康复者在研究期内不再返回易感状态；对已完成安全处置的遗体，不再计算后续接触传播。这是短期模型的简化。WHO 提醒，部分康复者体内仍可能存在病毒，故不能把 $R$ 的数学定义解释为所有康复者均不存在后续传播风险\cite{whofactsheet}。

\item \textbf{组内接触采用均匀混合近似。}
社区和医院接触按所在省份汇总，葬礼接触则用家庭型分组表示。这样可以在缺少逐人接触记录时建立人口层面的模型，但同一家庭型中的人不一定互有亲属关系。因此，家庭型分组只是对接触差异的近似，不能据此声称模型已经还原了真实亲属网络。

\item \textbf{遗体不参与通常的人口流动。}
遗体留在死亡所在地对应的分组，直到完成安葬；其余状态在基础设定中使用相同的省际流动规则。后一个设定主要出于缺少分状态出行数据的限制，并不意味着住院患者与普通居民具有相同的出行能力。解释空间预测时须考虑这一简化，尤其不能把由此得到的外溢估计视为已观测事实。

\item \textbf{个体病程用舱室平均转移过程表示。}
模型不逐人记录感染、发病、入院和结局的具体日期，而以相应状态的平均停留过程描述人群变化。这使常微分方程模型可以使用汇总病例资料，但会忽略病程的个体差异。参数的取值依据与不确定范围放在参数估计部分说明，不将拟合所得的平均值当作每位患者共同遵循的病程。

\item \textbf{报告病例与模型中的感染人数分开处理。}
报告记录可能受到漏检、延迟上报和历史补录的影响，不能直接等同于同期发生的全部感染。参数估计时设置观测关系，并区分不同层级和统计日期。病例曲线中的突变既可能来自疫情变化，也可能来自记录方式变化，不能仅靠拟合曲线判定原因。

\item \textbf{干预效果和输入情况按情景给定。}
干预强度允许随时间变化，不在整个研究期内固定为同一水平。进行方案比较时，除明确改变的条件外，其余设定保持一致；跨境输入另行计入，不把省际系统封闭的近似套用于输入病例分析。比较结果表示给定条件下的模型差异，不等同于已经测得的政策因果效应。

\item \textbf{疫苗与治疗只作为假设性情景扩展。}
扩展模型中的 $V$ 表示已获得保护者，不能直接解释为所有已经接种的人；$T$ 表示接受相应治疗的住院患者。将接种者全部视为立即获得保护，是该扩展中的理想化设定。WHO 指出，针对本迪布焦病毒的疫苗与特异性治疗仍在研发中\cite{whofactsheet}，因此不能把情景中预设的保护或治疗效果写成已经证实的临床结论。

\end{enumerate}

上述假设中，症状出现前的传播性质有流行病学依据；家庭型划分、组内混合和统一流动规则则是受数据条件限制的建模选择。
